\subsection{ISOMORPHISMS AND LOCAL ISOMORPHISMS}

In accordance with the general definitions (Set Theory, Chapter IV, § 1, no. 5) an isomorphism \(f\) of a topological group \(G\) onto a topological group \(G'\) is a bijective mapping of \(G\) onto \(G'\) which is simultaneously an isomorphism of the group structure of \(G\) onto that of \(G'\) , and a homeomorphism of \(G\) onto \(G'\) . In other words, \(f\) is an isomorphism of \(G\) onto \(G'\) if and only if: 1) \(f\) is bijective; 2) \(f(xy) = f(x)f(y)\) for all \(x, y \in G\) ; and 3) \(f\) is bicontinuous.

For example, if \(a\) is any point of \(G\) , the mapping \(x \to axa^{-1}\) is an isomorphism of \(G\) onto \(G\) , that is (loc. cit.), an automorphism of the topological group \(G\) . It is called an inner automorphism.

If a topology \(\mathfrak{C}\) is compatible with the group structure of a group \(\mathbf{G}\) , then \(\mathfrak{C}\) is also compatible with the group structure opposite to that of \(\mathbf{G}\) . If \(\mathbf{G}^0\) denotes the topological group obtained by giving the group opposite to \(\mathbf{G}\) the topology \(\mathfrak{C}\) , then the symmetry \(x \to x^{-1}\) is an isomorphism of the topological group \(\mathbf{G}\) onto the topological group \(\mathbf{G}^0\) .

DEFINITION 2. If G and G' are two topological groups, a local isomorphism of G with G' is a homeomorphism f of a neighbourhood V of the identity element of G onto a neighbourhood V' of the identity element of G' which satisfies the following conditions:

\begin{enumerate}
\def\labelenumi{\arabic{enumi})}
\tightlist
\item
  For each pair \(x, y\) of points of \(V\) such that \(xy \in V\) ,
\end{enumerate}

\[
f (x y) = f (x) f (y).
\]

\begin{enumerate}
\def\labelenumi{\arabic{enumi})}
\setcounter{enumi}{1}
\tightlist
\item
  If \(g\) is the mapping inverse to \(f\) , then for each pair of points \(x', y'\) of \(V'\) such that \(x'y' \in V'\) , we have \(g(x'y') = g(x')g(y')\) .
\end{enumerate}

The mapping g is then a local isomorphism of \(G'\) with G.

Two topological groups G, G' are said to be locally isomorphic if there exists a local isomorphism of G with G'.

Isomorphic topological groups are evidently locally isomorphic. The converse is false.

* For example, we shall see in Chapter V, § 1, that the topological groups R and T are locally isomorphic but not isomorphic. *

If \(f\) is a local isomorphism of \(G\) with \(G'\) , then every restriction of \(f\) to a neighbourhood of the identity element of \(G\) is again a local isomorphism of \(G\) with \(G'\) .

A local isomorphism of G with G is called a local automorphism of G.

In general, if \(f\) is a homeomorphism of a neighbourhood \(V\) of the identity element of \(G\) onto a neighbourhood \(V'\) of the identity element of \(G'\) which satisfies condition 1) of Definition 2, \(f\) does not necessarily satisfy condition 2) (see Exercise 7). Nevertheless, \(G\) and \(G'\) are in fact locally isomorphic.

PROPOSITION 3. Let G and G' be two topological groups, and let f be a homeomorphism of a neighbourhood V of the identity element of G onto a neighbourhood V' of the identity element of G', which satisfies condition 1) of Definition 2. Then f is an extension of a local isomorphism of G with G'.

For it is not hard to see that if W is a neighbourhood of the identity element of G such that \(W.W \subset V\) , then the restriction of f to W is a local isomorphism of G with \(G'\) .

